\section*{अध्याय \ref{ch:g9:our-responsibility} --- हमारी ज़िम्मेदारी}

\begin{solution}{exo:g9:our-responsibility:1}
रोज़ का इक़रारनामा अंगों के सही ठहराए हुए रखरखाव की तरह निभाना; उन
पदार्थों के जाल की क़ीमत लगाकर उसके दरवाज़े से इनकार करना; जनन वाला
नक़्शा और उसके पते इस्तेमाल करना; और परिवार के रुझानों को रखरखाव की
समय-सारणी की तरह पढ़ना।
\end{solution}

\begin{solution}{exo:g9:our-responsibility:2}
संचरण (\cref{ch:g9:microbes-and-infection}) --- मेरी \omterm{def:g1:healthy-habits:hygiene}{स्वच्छता} तुम्हारी
हिफ़ाज़त करती है; साझी \omterm{def:g9:immune-defenses:memory}{प्रतिरक्षा} (\cref{ch:g9:immune-defenses}) ---
मेरा टीका उनकी ढाल बनता है जिन्हें सिखाया नहीं जा सकता; और दवा की साझी
पूँजी (\cref{ch:g9:how-species-change} वाला रोध, और
\cref{ch:g9:unique-individuals} वाली \omterm{prop:g4:journey-of-food:absorb}{रक्त} की बही) --- जो अलग-अलग लोगों
के चुनावों से चलती रहती है, या भरी जाती है।
\end{solution}

\begin{solution}{exo:g9:our-responsibility:3}
क्योंकि विरासत में मिले रुझान जी जाने वाली सीमाएँ हैं, नियतियाँ नहीं: जो
परत जमने के रुझान वाला परिवार उसी हिसाब से खाता, चलता-फिरता और जाँच कराता
है वह नतीजा दोबारा लिख देता है --- यानी जानकारी उस बँटवारे को आदतों तथा
जाँचों की समय-सारणी में बदल देती है।
\end{solution}

\begin{solution}{exo:g9:our-responsibility:4}
मिसाल के लिए: \omterm{prop:g4:heart-and-blood:pulse}{नाड़ी} और दबाव --- यानी सड़क के किनारे पढ़ा गया रक्त-चक्र
(\cref{ch:g7:heart-and-circulation}); प्रतिवर्त वाला हथौड़ा --- यानी
तंत्रिका-परिपथ की जाँच (\cref{ch:g8:nervous-communication}); और \omterm{prop:g9:immune-defenses:vaccination}{टीकाकरण}
की बही --- यानी याददाश्त की चौकी, भरी हुई
(\cref{ch:g9:immune-defenses})।
\end{solution}

\begin{solution}{exo:g9:our-responsibility:5}
जाल का झटके सोखना बहुत-सी \omterm{def:g5:biodiversity:species}{प्रजातियों} पर चलता है --- यानी धागे और उनके
बदली वाले; और चयन का भविष्य विविधता पर --- यानी वह बचत, जिसमें से हर बदली
हुई परिस्थिति निकालती है। दोनों काम हैं, ख़र्च करने और बैंक में रखने
लायक़, कोई सजावट नहीं।
\end{solution}

\begin{solution}{exo:g9:our-responsibility:6}
नतीजा निकालने से पहले देखो; एक फ़र्क़ और एक \omterm{met:g4:what-plants-need:fairtest}{नियंत्रण} के साथ निष्पक्ष
परखो; किसी दावे को मानने से पहले उसकी कार्य-प्रणाली का नाम लो; आज़ाद
गवाहों को तौलो; और \omterm{def:g6:exploring-our-environment:components}{सजीव} मशीनरी का लिहाज़ करो, भीतर भी और बाहर भी।
\end{solution}

\begin{solution}{exo:g9:our-responsibility:7}
बुख़ार में घर रहा विद्यार्थी उस फैलाव का संचरण उसके स्रोत पर ही हटा देता
है --- यानी वह छुट्टी एक कटा हुआ रास्ता है, और बाद की हर दवा से सस्ता
(\cref{pb:g9:microbes-and-infection:1} वाला नारा, अमल में)।
\end{solution}

\begin{solution}{exo:g9:our-responsibility:8}
उस साझेपन की --- यानी हर सधा हुआ व्यक्ति एक बंद गली है, जो फैलाव को भूखा
मारती है। और टिके हुए: वह बहुत छोटा शिशु
(\cref{pb:g9:immune-defenses:1} वाला निगरानी में रखा संपर्क) और वह
मरीज़, जिसकी \omterm{def:g9:immune-defenses:memory}{प्रतिरक्षा} कमज़ोर है --- और जिन पर ढाल सिर्फ़ उनके पड़ोसियों
की सिखाई हुई है।
\end{solution}

\begin{solution}{exo:g9:our-responsibility:9}
``हर \omterm{prop:g3:sorting-living-things:groups}{कीट}'' में वे परागकर्ता भी आ जाते हैं जिन पर बाग़ का \omterm{prop:g3:flowers-fruits-seeds:transform}{फल} टिका है, और
वे कीट-भक्षक भी जो मुफ़्त काम करते थे
(\cref{ex:g5:biodiversity:orchard}); यह जाँच भविष्यवाणी करती है कि
\omterm{def:g5:flower-to-fruit:steps}{परागण} अब पैसे देकर कराना पड़ेगा और बदली वालों के मरने के बाद फैलाव और भी
बुरे होंगे --- यानी वह ``पूरी हिफ़ाज़त'', जो मज़दूरों को ही हटा देती है,
हिफ़ाज़त नहीं, टाला हुआ बिल है।
\end{solution}

\begin{solution}{exo:g9:our-responsibility:10}
हर अधूरी ख़ुराक उस वार्ड वाली चयन-छननी चला देती है --- कमज़ोरों को मारते
और रोधियों को आगे बढ़ाते हुए --- इसलिए हर ग़लत इस्तेमाल उस दवा की आगे की
उपयोगिता सबके लिए ख़र्च कर देता है। और भरने वाली आदतें: लिखे गए कोर्स
पूरे करो, और \omterm{def:g9:microbes-and-infection:sorted}{विषाणु} वाली बीमारियों के लिए कोई प्रतिजैविक मत माँगो।
\end{solution}

\begin{solution}{exo:g9:our-responsibility:11}
देखो: वह पोस्ट असल में दिखाती क्या है --- गवाहियाँ और रोशनियाँ, या
आँकड़े? परखो: क्या कोई निष्पक्ष परीक्षण हुआ --- एक फ़र्क़, एक \omterm{met:g4:what-plants-need:fairtest}{नियंत्रण}
--- या बस एक शुक्रगुज़ार मामला? कार्य-प्रणाली: वह इलाज किस कड़ी को रोकने
का दावा करता है --- क्या उसका नाम लिया जा सकता है? गवाह: क्या आज़ाद
स्रोत आपस में मिलते हैं, या सब कुछ बेचने वाले तक ही जाता है? लिहाज़: क्या
वह तुमसे चलता हुआ रखरखाव छोड़कर चमत्कार अपनाने को कहती है? पाँच क़दम ---
और ऐसी पोस्टें पहले दो से आगे कम ही बचती हैं।
\end{solution}

\begin{solution}{exo:g9:our-responsibility:12}
मिसाल के लिए: ``हर \omterm{def:g6:exploring-our-environment:components}{सजीव} \omterm{def:g5:first-look-at-cells:cell}{कोशिकाओं} से बना है, और उन सबमें एक ही पुराना
रसायन चलता है --- और तुम भी वैसे ही बने हो। तुम्हारा शरीर अंगों, सरहदों
और संदेशवाहकों की एक सँभाली हुई मशीन है --- और उसकी सँभाल तुम्हारी है।
तुम्हारा हर खाना और हर साँस तुम्हें \omterm{def:g5:food-webs:roles}{उत्पादकों}, \omterm{def:g5:food-webs:roles}{उपभोक्ताओं} तथा अपघटकों के
एक जाल से बाँधती है --- और उसकी सँभाल भी तुम्हारी ही है। और यह सब मिलकर
एक ही शाखाओं वाला परिवार है, एक ही कूट में लिखा हुआ, और चयन के नीचे
बदलता हुआ --- और तुम वही शाखा हो जिसने उस \omterm{def:g6:classification-kinship:tree}{वृक्ष} को समझ लिया।''
\end{solution}

\begin{solution}{pb:g9:our-responsibility:1}
\textbf{1.} बूँदें और हाथ --- इसलिए खाँसी ढको, धोओ, कमरों में हवा आने
दो, और बीमार हो तो घर रहो। और कर्मचारी पहले, क्योंकि वही उस इमारत के सबसे
नाज़ुक लोगों को जोड़ते हैं: उनकी सिखाई ठीक वहीं साझी होती है जहाँ न सिखाए
जा सकने वाले इकट्ठा हैं।

\textbf{2.} वह बही दुर्घटनाओं और शल्य-चिकित्साओं के दिन के लिए मेल खाता
\omterm{prop:g4:journey-of-food:absorb}{रक्त} भरती है; और वह \omterm{prop:g9:immune-defenses:vaccination}{टीकाकरण} वाली बही सधी हुई याददाश्तें। और हक़ सबका ---
बही का मेल समूह से मिलता है, और उस साझेपन की ढाल बच्चों, बूढ़ों तथा
बीमारों को ढकती है।

\textbf{3.} मिसाल के लिए: ``एक वार्ड ने कभी सिर्फ़ धुले हाथों से अपनी
माँओं की मौतें आधी कर दी थीं: यह `हो-हल्ला' उन हरकारों को रोकता है
जिन्हें कोई \omterm{def:g1:the-five-senses:sense}{आँख} नहीं देख सकती, और ज़ेमेलवाइस के बाद की डेढ़ सदी ने बस यही
पक्का किया है कि दुनिया की सबसे सस्ती दवा साबुनदानी से निकलती है।''

\textbf{4.} प्रतिजैविकों वाली साझी पूँजी, और वह भी अधूरे चयन से: हर अधूरा
कोर्स उस रोधी अल्पसंख्या को बख़्शता और आगे बढ़ाता है, और इस तरह उस दवा
की ताक़त सबके आगे के \omterm{def:g9:microbes-and-infection:stages}{संक्रमणों} के ख़िलाफ़ ख़र्च कर देता है।

\textbf{5.} वह छिड़काव कीड़ों के साथ-साथ परागकर्ताओं और कीट-भक्षकों को भी
ले जाएगा: मधुमक्खियों के साथ \omterm{prop:g3:flowers-fruits-seeds:transform}{फल} का बैठना गिर जाता है, और बचे हुओं का अगला
फैलाव न कोई सोनकीड़ा पाता है न कोई चिड़िया --- यानी मुफ़्त सेवाएँ रुक
जाती हैं, और बिल किराए के छत्तों तथा और छिड़काव की शक्ल में लौट आता है।

\textbf{6.} वह मटर जैसी हरी बढ़वार गरम, ठहरे पानी में पूरी रात साँस लेती
है, और उस पानी में \omterm{def:g7:what-is-respiration:respiration}{ऑक्सीजन} वैसे भी कम होती है: भोर वाला न्यूनतम तल तक जा
पहुँचा (\cref{ex:g7:oxygen-in-the-water:curve})। नुस्ख़ा: रात भर चलता
फ़व्वारा, कीचड़ तथा घास-फूस पतला करना --- और उन पोषक तत्वों के आने पर भी
नज़र, जिन्होंने वह बढ़वार खिलाई।

\textbf{7.} \omterm{def:g2:where-animals-live:habitat}{आवास}: वह दलदल प्रजनन की ज़मीन है --- \omterm{def:g2:where-animals-live:habitat}{आवास} नहीं, तो बाशिंदे
नहीं। जाल: उसकी \omterm{def:g5:biodiversity:species}{प्रजातियाँ} आस-पास की शृंखलाओं का \omterm{def:g5:flower-to-fruit:steps}{परागण} करती हैं, उन्हें
साफ़ करती हैं और खिलाती हैं। बचत: उसकी विविध \omterm{def:g8:reproduction-and-environment:population}{आबादियाँ} वे \omterm{def:g9:genes-and-dna:gene}{विकल्पियाँ} बैंक
में रखती हैं जिन्हें बुरे वक़्त बुलाएँगे। (और जहाँ चरागाहें लौटती हैं
वहाँ सारस भी लौटते हैं।)

\textbf{8.} वही बिना चुकाई बही: बाढ़ का पानी थामकर धीरे-धीरे छोड़ना, पानी
छानना, परागकर्ताओं को घर देना, और \omterm{prop:g3:sorting-living-things:groups}{मछलियों} के बच्चों को पालना --- ये सेवाएँ
बिल तभी भेजती हैं जब वह दलदल सुखा दिया जाए और क़स्बा उन्हें इंजीनियरों से
ख़रीदे।

\textbf{9.} जवाब खुले हैं; एक बचाव लायक़ जोड़ी: सबसे अच्छा --- वह
शरणस्थली (टिकाऊ इस्तेमाल के तीनों नियमों के तहत बैंक में रखी पूँजी); और
सबसे बुरा --- अधूरे कोर्स या वह छिड़काव वाला प्रस्ताव (साझी पूँजी ख़र्च,
और मज़दूर बेदख़ल), और यह फ़ैसला भरपाई-से-तेज़-मत-लो,
मज़दूरों-को-बचाए-रखो, और जो-उधार-लिया-वह-लौटाओ के हिसाब से।

\textbf{10.} देखा: भोर में मरी \omterm{prop:g3:sorting-living-things:groups}{मछलियाँ}, हरा पानी, और नीची पढ़तें।
निष्पक्ष परीक्षण: भोर बनाम दोपहर बाद की पढ़तें; और यह तालाब बनाम कोई साफ़
तालाब। कार्य-प्रणाली का नाम: गरम पानी की गिरी हुई छत के नीचे उस बढ़वार का
रात वाला \omterm{def:g7:what-is-respiration:respiration}{श्वसन}। गवाह: मापक यंत्र के आँकड़े, गिनती में बचे सहनशील, और वह
समय --- सब आपस में मिलते हुए। लिहाज़: वह नुस्ख़ा \omterm{def:g9:heredity-and-traits:traits}{लक्षणों} को दवा देने के
बजाय तालाब की \omterm{def:g6:exploring-our-environment:components}{सजीव} मशीनरी सँभालता है।

\textbf{11.} मिसाल के लिए: ``हमारी फ़्लू की लहर वहीं ख़त्म हुई जहाँ हाथ
धुले और बीमारी की छुट्टियाँ ली गईं --- सेहत एक साझा खाता है। जो दलदल हमने
बचाया वह पूँजी है, नज़ारा नहीं, और वह थामी हुई बाढ़ों तथा परागित \omterm{def:g6:producing-our-food:farming}{फ़सलों}
में चुकेगा। साल के सबसे ऊँचे दावे एक ही सवाल पर सबसे बुरे निकले:
कार्य-प्रणाली का नाम बताओ। तरीक़ा सँभाले रखो: यही इस क़स्बे का सबसे सस्ता
ढाँचा है।''

\textbf{12.} जवाब खुला है --- एक ईमानदार संभावना: ``वह घोंघा और तुम,
दोनों उन्हीं चार अक्षरों से बने हो: पूरे परिवार का ख़्याल रखना।''
\end{solution}
